\phant{GIỚI THIỆU}\label{sec:gioithieu}

\subsection{Bối cảnh}

Tính toán hiệu năng cao (HPC) được xây dựng trên hai trụ cột lập trình song song:
bộ nhớ chia sẻ (shared memory) và bộ nhớ phân tán (distributed memory). Trong mô
hình thứ nhất, nhiều luồng (thread) cùng truy cập một không gian địa chỉ, phối
hợp với nhau qua biến dùng chung, rào chắn và cơ chế đồng bộ. Trong mô hình thứ
hai, nhiều tiến trình (process) chạy trên các nút khác nhau, trao đổi dữ liệu
bằng thông điệp. Mỗi mô hình có chi phí riêng: đồng bộ và tranh chấp bộ nhớ ở
mô hình chia sẻ; độ trễ và băng thông truyền thông ở mô hình phân tán. Việc viết
mã song song đúng, nhanh và mở rộng được đòi hỏi hiểu biết đồng thời về thuật
toán, kiến trúc máy tính, mô hình bộ nhớ và công cụ biên dịch.

Trong hơn một thập kỷ, chuẩn C++ đã bổ sung dần các cơ chế song song ngay trong
ngôn ngữ: \code{std::thread} và \code{std::async} (C++11), thuật toán song song
với execution policy (C++17), coroutine cùng ranges, semaphore, latch, barrier
(C++20) \cite{diehl2023shared,diehl2024book}. Song song với đó, OpenMP và MPI vẫn
là hai chuẩn thực tế được dùng rộng rãi nhất trong các ứng dụng khoa học. Sự
phong phú này tạo ra nhiều cách diễn đạt cùng một phép song song, nhưng cũng làm
tăng không gian lỗi: một chương trình có thể biên dịch thành công, chạy không
báo lỗi, mà vẫn cho kết quả sai hoặc không hề tận dụng số lõi được yêu cầu.

Cùng lúc đó, các mô hình ngôn ngữ lớn (LLM) đã trở thành công cụ hỗ trợ sinh mã
phổ biến. Từ một mô tả bằng ngôn ngữ tự nhiên, mô hình có thể tạo ra hàm, lớp,
thậm chí cả chương trình hoàn chỉnh. Các đánh giá ban đầu tập trung vào bài toán
lập trình thuật toán hoặc mã tuần tự, nơi tiêu chí đúng/sai tương đối rõ ràng.
Với mã HPC, tiêu chí phức tạp hơn nhiều: ngoài tính đúng, còn có khả năng biên
dịch trên toolchain thực tế, khả năng kiểm soát tài nguyên, tính ổn định khi
tăng số lõi và hiệu năng so với mã do con người viết.

\subsection{Vấn đề nghiên cứu}

Khoảng trống mà bài báo gốc \cite{diehl2026} hướng tới nằm ở giao điểm của ba yếu
tố. Thứ nhất, phần lớn nghiên cứu về LLM sinh mã chỉ đánh giá một mô hình lập
trình đơn lẻ, chẳng hạn OpenMP, OpenACC hoặc MPI, nên thiếu so sánh ngang hàng
giữa các mức trừu tượng. Thứ hai, các nghiên cứu thường dừng ở tỉ lệ biên dịch
hoặc tỉ lệ vượt qua bộ kiểm thử, mà không đo khả năng mở rộng trên phần cứng
song song thật. Thứ ba, kiến thức của LLM về hệ sinh thái công cụ — trình biên
dịch nào, thư viện nền nào, phiên bản chuẩn nào — hầu như chưa được kiểm chứng,
dù đây chính là nguồn lỗi phổ biến trong thực tế.

Bài toán Mandelbrot được chọn làm phép thử có kiểm soát. Đây là bài toán song
song hiển nhiên (\term{embarrassingly parallel}): mỗi điểm ảnh được tính độc lập,
không cần bộ giải, không có trao đổi dữ liệu giữa các đơn vị tính toán. Nhờ vậy,
sai sót quan sát được chủ yếu đến từ cách biểu diễn song song, cách quản lý luồng
và cách xuất kết quả, chứ không bị nhiễu bởi thuật toán số phức tạp. Nói cách
khác, bài toán đóng vai trò một \term{benchmark tối thiểu} để cô lập năng lực sinh
mã song song của LLM.

\subsection{Mục tiêu và câu hỏi nghiên cứu}

Mục tiêu của bài báo là đánh giá một cách có hệ thống bốn LLM hàng đầu
(ChatGPT-4, ChatGPT-5, Claude và LLaMA \cite{achiam2023gpt4,touvron2023llama}) trên năm biến thể lập trình song song:
C++11 bất đồng bộ, C++17 thuật toán song song, C++20 coroutine, OpenMP và
MPI+OpenMP. Ba câu hỏi nghiên cứu được đặt ra:

\begin{enumerate}[label=\textbf{RQ\arabic*.},leftmargin=2.2cm]
  \item LLM sinh mã C++ song song đúng cú pháp và biên dịch được ở mức nào?
  \item Những lỗi thực thi và ngữ nghĩa nào xuất hiện, chúng khó sửa đến đâu, và
        nguyên nhân gốc là gì?
  \item Sau khi sửa, mã sinh ra có mở rộng theo số lõi và số nút không, và có
        cạnh tranh được với mã do con người viết không?
\end{enumerate}

\subsection{Đóng góp}

Bài báo gốc nêu bốn đóng góp: (i) thiết kế prompt có cấu trúc cho C++11/17/20,
OpenMP và MPI+OpenMP, dùng nhất quán cho mọi mô hình; (ii) phân tích chi tiết kết
quả biên dịch, lỗi thực thi và tính đúng chức năng; (iii) đánh giá coroutine
C++20 — cơ chế song song mới nhất, đang được tích hợp vào Kokkos và HPX; (iv)
phân loại lỗi và cách khắc phục, qua đó chỉ ra giới hạn của LLM trong tổng hợp và
gỡ lỗi mã HPC.

Báo cáo này bổ sung thêm ba đóng góp trình bày:

\begin{itemize}
  \item Hệ thống hóa cơ sở lý thuyết của từng mô hình lập trình, giải thích vì sao
        các lỗi quan sát được lại xảy ra ở đúng điểm đó, thay vì chỉ liệt kê triệu
        chứng.
  \item Xây dựng một phân loại lỗi theo tầng (cú pháp, API, ngữ nghĩa, đồng bộ,
        định dạng, đo đạc) kèm mức độ khó sửa và kiến thức cần có.
  \item Chuyển kết quả nghiên cứu thành quy trình kiểm duyệt và danh mục kiểm tra
        thực dụng cho nhóm phát triển HPC có sử dụng LLM hỗ trợ.
\end{itemize}

\subsection{Phạm vi và cách tiếp cận}

Phạm vi thực nghiệm giới hạn ở bốn mô hình ngôn ngữ, năm biến thể lập trình, một
bài toán, một trình biên dịch (GCC 11.5.0), một backend song song (Intel TBB
2022.2) và một cụm CPU. Mỗi biến thể được sinh một lần từ một prompt cố định,
sau đó con người sửa lỗi biên dịch và lỗi thực thi theo cách tối thiểu, ghi nhận
lại từng thay đổi. Kết quả hiệu năng vì thế phản ánh mã \emph{sau khi sửa}, không
phải chất lượng nguyên bản của lần sinh đầu tiên. Cách tiếp cận này phù hợp với
câu hỏi thực tiễn: điều người dùng quan tâm là LLM tạo ra bao nhiêu việc phải sửa
và sau khi sửa thì mã có dùng được không.

\subsection{Bố cục báo cáo}

Phần~\ref{sec:coso} trình bày cơ sở lý thuyết: bài toán Mandelbrot, bốn cơ chế
song song của C++, OpenMP, MPI và mô hình lai, cùng nền tảng về LLM sinh mã.
Phần~\ref{sec:phuongphap} mô tả thiết kế thực nghiệm, prompt, quy trình kiểm định,
độ đo và môi trường đo. Phần~\ref{sec:ketqua} trình bày kết quả theo từng mô hình
và từng biến thể, kèm phân tích lỗi, chất lượng mã và khả năng mở rộng.
Phần~\ref{sec:thaoluan} thảo luận nguyên nhân, hàm ý thực tiễn và hạn chế.
Phần~\ref{sec:ketluan} kết luận và nêu hướng phát triển. Phụ lục cung cấp bảng
tổng hợp lỗi, trích mã tiêu biểu và danh mục kiểm tra.
